%========================================================================================
\section{Validation, transactions et intégrité}
\sectionmark{Validation et intégrité}
\label{sec:validation}
%========================================================================================

\subsection{Trois niveaux de validation}

Avant d'être écrite, une donnée passe trois contrôles l'un après l'autre (tableau~\ref{tab:niveaux}). Chacun a son utilité, et aucun ne suffirait tout seul.

\begin{table}[H]
  \centering
  \small
  \caption{Niveaux de validation.}
  \label{tab:niveaux}
  \begin{tabularx}{\linewidth}{@{}l X X@{}}
    \toprule
    \textbf{Niveau} & \textbf{Règles} & \textbf{Pourquoi} \\
    \midrule
    Serializer & Types, longueurs, choix, champs en lecture seule, cohérence des dates en \code{PATCH} & Refuser tôt, avec un message par champ \\
    \code{Model.clean()} & Destination dans les dates du voyage, activité dans les dates de la destination, pas de raccourcissement qui laisserait un enfant en dehors & Règles métier qui ont besoin de lire d'autres lignes \\
    Contraintes PostgreSQL & \code{CHECK} et \code{UNIQUE} du tableau~\ref{tab:contraintes} & Dernière protection, qui marche aussi quand deux requêtes arrivent en même temps \\
    \bottomrule
  \end{tabularx}
\end{table}

En \code{PATCH}, le client peut n'envoyer qu'une des deux dates. Le serializer reprend alors l'autre date dans la valeur enregistrée avant de les comparer. Sinon, une date de fin plus tôt que la date de début existante passerait la validation et échouerait seulement dans la base.

\subsection{Appliquer \texttt{Model.clean()} depuis DRF}

DRF n'appelle ni \code{Model.clean()} ni \code{validate_constraints()}. Sans rien faire, l'API ignorerait donc les règles de \code{clean()}, et une contrainte violée donnerait une \code{IntegrityError} au moment de l'écriture. Nous avons écrit le mixin \code{ModelCleanMixin} pour régler ce problème.

\begin{lstlisting}[language=Python,caption={Validation par le modèle (\texttt{apps/core/serializers.py}, extrait).},label={lst:modelclean}]
class ModelCleanMixin:
    def validate(self, attrs):
        attrs = super().validate(attrs)
        if self.instance is not None:
            candidate = copy.copy(self.instance)
        else:
            candidate = self.Meta.model()
        for field, value in attrs.items():
            setattr(candidate, field, value)
        try:
            candidate.clean()
            candidate.validate_constraints()
        except DjangoValidationError as exc:
            raise serializers.ValidationError(django_error_detail(exc)) from exc
        return attrs
\end{lstlisting}

Les règles sont testées sur une \emph{copie} de l'instance. Si la validation échoue, l'objet chargé par la vue ne change pas. Les messages de Django sont remis au format de DRF, et la clé \code{__all__} devient \code{non_field_errors}.

\subsection{Contraintes violées malgré la validation}

Deux requêtes envoyées en même temps peuvent toutes les deux passer la validation, puis entrer en conflit au moment d'écrire. Par exemple, deux éditeurs mettent chacun une destination en position 2. La contrainte \code{itinerary_destination_unique_trip_order} refuse alors la deuxième écriture. Le gestionnaire d'exceptions transforme l'\code{IntegrityError} en réponse 400. Il retrouve le bon message grâce au nom de la contrainte que PostgreSQL renvoie dans son diagnostic.

\begin{lstlisting}[language=Python,caption={Message de la contrainte violée (\texttt{apps/core/exceptions.py}, extrait).},label={lst:integrity}]
def integrity_error_message(exc):
    """Retrouve le message de la contrainte violée à partir du nom renvoyé par PostgreSQL."""
    diag = getattr(exc.__cause__, "diag", None)  # psycopg expose le diagnostic du serveur
    name = getattr(diag, "constraint_name", None)
    return _constraint_messages().get(name, INTEGRITY_MESSAGE)
\end{lstlisting}

Le client reçoit donc \og{}Une destination occupe déjà cette position dans le voyage.\fg{} avec le code \code{integrity_error}, au lieu d'une erreur 500.

\subsection{Transactions}

Les opérations qui écrivent plusieurs lignes sont atomiques. Les services \code{create_trip}, \code{clone_trip} et \code{change_participant_role} ont le décorateur \code{@transaction.atomic}, comme chaque mutation GraphQL. Le service \code{add_participant} ouvre un point de sauvegarde (un \code{transaction.atomic()} imbriqué) autour de l'insertion. Si l'unicité est violée, on attrape l'erreur et la transaction englobante reste utilisable.

\subsection{Démonstration du rollback}

Quand on clone un voyage, on écrit dans quatre tables : le voyage, le participant, les destinations et les activités. Le test de l'extrait~\ref{lst:rollback} provoque une panne à la toute fin, pendant la copie des activités. Il compare ensuite le contenu des quatre tables avant et après la requête.

\begin{lstlisting}[language=Python,caption={Panne simulée pendant le clonage (\texttt{test\_trip\_actions.py}, extrait).},label={lst:rollback}]
def test_failure_midway_gives_500_and_leaves_database_unchanged(
    self, client_for, owner, source, monkeypatch
):
    def snapshot():
        return [
            list(model.objects.order_by("pk").values())
            for model in (Trip, TripParticipant, Destination, Activity)
        ]

    def failing_bulk_create(objs, *args, **kwargs):
        raise RuntimeError("panne simulée pendant la copie des activités")

    before = snapshot()
    monkeypatch.setattr(Activity.objects, "bulk_create", failing_bulk_create)
    client = client_for(owner)
    client.raise_request_exception = False

    response = client.post(clone_url(source))

    assert response.status_code == 500
    assert response.json() == {"error": INTERNAL_ERROR}
    assert snapshot() == before
\end{lstlisting}

Quand la panne arrive, le nouveau voyage, son participant et ses destinations sont déjà insérés. La transaction annule tout. Le client reçoit une erreur 500 en JSON, sans trace de pile, et la base revient exactement à son état de départ. Nous vérifions le même scénario pour la mutation GraphQL \code{cloneTrip} (\code{apps/trips/tests/test_graphql.py}) et pour le service seul (\code{test_services.py}).
